\chapter{Human AI Interaction}

AI systems are becoming increasingly prevalent in our daily lives, from virtual assistants to recommendation systems. As these systems become more sophisticated, it is crucial to understand how humans interact with them and how to design interfaces that facilitate effective communication and collaboration between humans and AI.

\section{What is Artificial Intelligence?}

Artificial Intelligence (AI) refers to the simulation of human intelligence in machines that are programmed to think and learn. AI systems can perform tasks that typically require human intelligence, such as visual perception, speech recognition, decision-making, and language translation. Even though intelligence in humans can be \emph{measured}, AI can't be measured. There are however various tests that have been proposed to evaluate if a machine can be considered intelligent. The importance of these tests is shown in the Chinese Room thought experiment.
\begin{itemize}
    \item \ul{Chinese Room Experiment}:
    
    The Chinese Room thought experiment involves a person who does not understand Chinese sitting in a room with a set of instructions for how to respond to Chinese characters. The person receives Chinese characters as input and uses the instructions to produce appropriate responses, which are then sent back out of the room.
    
    From the outside, it may appear that the person in the room understands Chinese, but in reality, they are simply following a set of rules without any comprehension of the language. This thought experiment raises questions about whether a machine can truly understand language or if it is simply following a set of instructions.
\end{itemize}

The most well-known test for evaluating machine intelligence is the Turing Test, proposed by Alan Turing:
\begin{itemize}
    \item \ul{Turing Test}:
    
    In the Turing Test, a human evaluator interacts with both a machine and a human without knowing which is which. If the evaluator cannot reliably distinguish between the machine and the human, then the machine is said to have passed the Turing Test and is considered to be intelligent.

    It only tests for the ability to mimic human behavior, not for true understanding or consciousness.
\end{itemize}

\section{User's control}

When designing AI systems, it is important to consider the user's control over the system. AI's goal is usually translating the raw data into a model that can be used to make predictions or decisions. However, there are different levels of control that users can have over this process:
\begin{itemize}
    \item An algorithm that is fully automated does this conversion without any user input, so no human control is involved.
    \item Humans provide feedback to the training algorithm, which can be used to improve the model's performance. This is often referred to as ``human-in-the-loop'' learning.
    \item Users replace the training data or algorithms with their own knowledge and expertise, effectively taking all the control of the learning process, but having to do all the work.
\end{itemize}

Some researches argue that there is a \emph{trade-off between the level of control and the automation of the system}. Fully automated systems may be more efficient and require less effort from users, but they may also be less transparent and less customizable. On the other hand, systems that allow for more user control may be more flexible and adaptable to individual needs, but they may also require more effort and expertise from users. However, some researchers argue that this trade-off is not necessarily inherent in the Human-AI interaction. This idea is shown in Figure~\ref{fig:control-automation}, where the diagonal line represents the trade-off between control and automation, but there are also areas where both can be achieved simultaneously.
\begin{figure}
    \centering
    \begin{tikzpicture}[
        font=\sffamily,
        % Definición de colores basados en la imagen original
        baseblue/.style={fill={rgb,255:red,79; green,142; blue,205}},
        boxborder/.style={draw=black, fill=none, line width=1pt},
        labelstyle/.style={text=black, font=\sffamily\small},
        titlewhite/.style={text=black, font=\sffamily\small},
        titlelight/.style={text=black, font=\sffamily\small}
    ]
        % Eje X e Y
        \draw[-Stealth, line width=1.5pt] (0,0) -- (12.5,0) node[anchor=north] {};
        \draw[-Stealth, line width=1.5pt] (0,0) -- (0,6) node[anchor=east] {};



        % 6. Flecha diagonal amarilla (Atraviesa de abajo-derecha a arriba-izquierda)
        \draw[{Triangle[length=4mm, width=5mm]}-{Triangle[length=4mm, width=5mm]}, draw={rgb,255:red,253; green,205; blue,93}, line width=4.5pt] 
            (0.2,5.6) -- (11.5,0.4);

        % 2. Recuadros internos blancos (Bordes finos de los cuadrantes)
        % Cuadrante Superior Izquierdo
        \draw[boxborder] (0.2,3.1) rectangle (5.8,5.8);
        % Cuadrante Superior Derecho
        \draw[boxborder] (6.2,3.1) rectangle (11.8,5.8);
        % Cuadrante Inferior Derecho
        \draw[boxborder] (6.2,0.2) rectangle (11.8,2.9);

        % 3. Textos del Cuadrante Superior Izquierdo (Human Mastery)
        \node[titlelight, anchor=north west] at (2.8,5.7) {Human Mastery};
        \node[labelstyle, anchor=north west, align=left, text width=3cm, row sep=1pt] at (0.3,4.7) {
            Bicycle \\ Piano \\ Chess
        };

        % 4. Textos del Cuadrante Superior Derecho (Reliable, Safe...)
        \node[titlelight, anchor=north east] at (11.2,5.7) {Reliable, Safe \& Trustworthy};
        \node[labelstyle, anchor=north east, align=right, text width=4cm] at (11.7,5.2) {
            Elevator \\ Camera \\ Plane Autopilot \\ Calculator
        };

        % 5. Textos del Cuadrante Inferior Derecho (Computer Control)
        \node[labelstyle, anchor=south east, align=right, text width=4cm] at (11.7,1.8) {
            Car Brake assistant \\ Airbag
        };
        \node[titlelight, anchor=south east] at (9.4,0.3) {Computer Control};

        % 7. Ejes y Etiquetas externas (Control y Automation)
        % Eje Vertical: Control
        \node[labelstyle, anchor=east] at (-0.2, 5.8) {human};
        \node[labelstyle, font=\sffamily\bfseries, anchor=east] at (-0.2, 3.0) {Control};
        \node[labelstyle, anchor=east] at (-0.2, 0.2) {AI};

        % Eje Horizontal: Automation
        \node[labelstyle, anchor=north] at (0.0, -0.1) {low};
        \node[labelstyle, font=\sffamily\bfseries, anchor=north] at (6.0, -0.1) {Automation};
        \node[labelstyle, anchor=north] at (12.0, -0.1) {high};

    \end{tikzpicture}
    \caption{Control vs Automation in Human-AI Interaction}
    \label{fig:control-automation}
\end{figure}

\section{Explainability in AI Systems}

Explainability refers to the ability of an AI system to provide clear and understandable explanations for its decisions and actions. There are several reasons why explainability is important in AI systems:
\begin{itemize}
    \item \ul{Trust}: Users are more likely to trust an AI system if they can understand how it works and why it makes certain decisions.
    \item \ul{Accountability}: Explainable AI systems allow decisions to be audited and held accountable, which is especially important in high-stakes applications such as healthcare and finance.
    \item \ul{Ethical Considerations}: Explainability can help ensure that AI systems are used ethically and responsibly, as it allows users to understand the potential consequences of the system's actions.
    \item \ul{Debugging and Improvement}: Explainable AI systems can provide insights into how the system is functioning, which can be useful for debugging and improving the system over time.
\end{itemize}

There are different recipients of explanations in AI systems, including:
\begin{itemize}
    \item \ul{Users affected by the model's decisions}: These users may require explanations to understand how the model's decisions impact them and to ensure that they are treated fairly.
    \item \ul{Data scientists and developers}: These individuals may require explanations to understand how the model is functioning and to identify potential issues or areas for improvement.
    \item \ul{Regulators and policymakers}: These stakeholders may require explanations to ensure that AI systems are compliant with regulations and ethical standards.
    \item \ul{Domain experts or users of the system}: These users may require explanations to understand how the AI system is making decisions in their specific domain and to ensure that it is providing accurate and relevant information.
    \item \ul{Managers and decision-makers}: These individuals may require explanations to understand the implications of the AI system's decisions and to make informed decisions based on the system's outputs. 
\end{itemize}

There is however an important trade-off between the accuracy of an AI system and its explainability. Highly accurate AI systems may be less transparent and harder for users to understand, which can lead to a lack of trust. On the other hand, more explainable AI systems may be less accurate but can foster greater trust among users. This is the goal of the field of Explainable AI (XAI), which aims to develop AI systems that are both accurate and explainable. Depeding on the application and the users, there are different types of explanations that may be more appropriate, such as:
\begin{itemize}
    \item \ul{Analytical Explanations}: These explanations provide a detailed analysis of the model's decision-making process, including the features and data that influenced the decision.
    \item \ul{Visual Explanations}: These explanations use visualizations to help users understand the model's decision-making process, such as heatmaps or feature importance graphs.
    \item \ul{Rule-Based Explanations}: These explanations provide a set of rules or guidelines that the model follows to make decisions, which can be easier for users to understand. It is presented as a simple \texttt{if-then} statement. It may be too simplistic for complex models, but it can be useful for certain applications where interpretability is more important than accuracy.
    \item \ul{Textual Explanations}: These explanations provide a natural language description of the model's decision-making process, which can be more accessible for users who are not familiar with technical details. There are different types of textual explanations, such as:
    \begin{itemize}
        \item \ul{Contrastive Explanations}: These explanations focus on why a particular decision was made instead of another, highlighting the differences between the two outcomes.
        \item \ul{Counterfactual Explanations}: These explanations focus on what changes would need to be made to the input data to produce a different outcome, providing insights into the model's decision boundaries.
    \end{itemize}
\end{itemize}

There are four steps that should be followed to design effective explanations for AI systems:
\begin{itemize}
    \item \ul{Identify the target audience}: Understanding who will be receiving the explanations is crucial for designing explanations that are appropriate and effective for that audience.
    \item \ul{Determine the purpose of the explanation}: The purpose of the explanation will guide the content and format of the explanation, whether it is to build trust, provide accountability, or facilitate debugging and improvement.
    \item \ul{Choose the content of the explanation}: The content of the explanation should be relevant and informative for the target audience, providing insights into the model's decision-making process without overwhelming users with technical details.
    \item \ul{Select the format of the explanation}: The format of the explanation should be accessible and understandable for the target audience, whether it is through visualizations, textual descriptions, or rule-based explanations.
\end{itemize}

A last important consideration when designing explanations for AI systems is reliance on the AI's answers.
\begin{itemize}
    \item \ul{Over-reliance}: Users may become overly reliant on the AI system's outputs, leading to a lack of critical thinking and potential errors if the system provides incorrect information.
    \item \ul{Under-reliance}: Conversely, users may under-rely on the AI system's outputs, leading to missed opportunities for improved decision-making and efficiency.
\end{itemize}

Appropriate explanations can help mitigate these issues by providing users with the necessary information to make informed decisions about when to trust the AI system and when to exercise caution.